\section{Methodology}
\label{sec:methodology}

\subsection{DEM model and specimen preparation}
\label{subsec:dem_model}

% PFC3D, rolling resistance contact model, periodic boundaries
% Particle size distribution (1.0--3.0 mm, Toyoura sand analogue)
% Ribbed shear walls for applying shear force
% Isotropic consolidation to p' = 100 kPa
% Undrained (constant-volume) condition via servo mechanism
% DEM parameters table (shared with K0 paper)

\subsection{Loading path design}
\label{subsec:loading_paths}

% Core methodological contribution: systematic elimination of confounding factors

\paragraph{Single-8 shear stress path.}
% Eqs for tau_8x, tau_8y
% Figure-8 pattern in tau_x--tau_y space

\paragraph{Equivalent-magnitude unidirectional path.}
% Key innovation: tau_uni = sign(tau_8x) * sqrt(tau_8x^2 + tau_8y^2)
% Ensures identical shear stress magnitude at every instant
% Only difference: direction is fixed

\paragraph{Double-8 shear stress path.}
% Swaps x and y components in alternating cycles
% Preserves: (1) magnitude, (2) rate of change
% Only difference vs single-8: loading history (axis alternation)

\paragraph{Summary of controlled variables.}
% Table or diagram summarizing what is controlled between each pair of comparisons
% Uni vs single-8: same magnitude, different direction and rate
% Single-8 vs double-8: same magnitude and rate, different directional history

\subsection{Parameter calibration and validation}
\label{subsec:calibration_validation}

% Validation against Nakata et al. (1998) monotonic HCA data
% Clarify: validation targets contact model parameters, not boundary conditions
% Parameters inherited from K0 paper (cross-reference)

\subsection{Microscopic descriptors}
\label{subsec:micro_descriptors}

% Zm: mechanical coordination number (same definition as K0 paper)
% Fc: invariant of anisotropic fabric tensor (unsigned scalar -- appropriate here
%     because the question is degree of anisotropy, not its direction)
% rho_c: contact density distribution
